\documentclass[12pt,a4paper]{article}

% 支持中文
\usepackage[UTF8]{ctex}

% 页面设置
\usepackage[top=2.5cm, bottom=2.5cm, left=3cm, right=3cm]{geometry}

% 常用宏包
\usepackage{amsmath, amssymb}
\usepackage{amsfonts}
\usepackage{algorithm}
\usepackage{algpseudocode}
\usepackage{mathrsfs}
\usepackage{graphicx}
\usepackage{subcaption}
\usepackage[colorlinks=true,linkcolor=blue]{hyperref}
\usepackage{indentfirst} % 使章节后的首段也缩进
\usepackage{setspace}     % 行距设置
\usepackage{listings}     % 插入代码块
\usepackage{xcolor}       % 设置代码颜色
\usepackage{mathtools}
\usepackage{bm}
\usepackage{multirow}
\usepackage{siunitx}
\usepackage{caption}
\setlength{\parindent}{2em} % 设置段落首行缩进为 2em

\usepackage[backend=biber,        % 必须用 biber 后端
            style=gb7714-2015]    % 国标顺序编码制
            {biblatex}  
\addbibresource{refs.bib}

\usepackage{booktabs} % 用于三线表 toprule, midrule, bottomrule
\usepackage{tabularx} % 用于自适应宽度的表格
\usepackage{array}    % 用于表格排版
\usepackage{makecell} % 用于单元格内强制换行

\usepackage{tikz}
\usetikzlibrary{shapes.geometric, arrows.meta, positioning, calc}

\usepackage{enumitem}
\usepackage{adjustbox}

\begin{document}



\section{科学问题}
\subsection{问题解析}
本文研究在部分可观测、对抗性的网络防御场景中，如何将自主防御建模为多目标序贯决策问题，并高效发现与利用能够适应不同运维偏好的 Pareto 防御策略集。

可以拆分成三个子问题

子问题 1：网络防御为什么不能视为单目标优化

真实蓝方防御不是只追求“拦住攻击”这一件事，还必须同时考虑：

安全效果
业务扰动
防御开销

所以科学问题的第一层是：

如何把自主网络防御建模为一个多目标而非单目标的决策问题。

子问题 2：在对抗与部分可观测条件下，怎样发现一组而不是一个防御策略

在 CybORG++ 这类环境里，防御者面对的是：

红方攻击行为
状态不完全可观测
防御偏好可能变化

所以第二层科学问题是：

在部分可观测、对抗性的 cyber environment 中，如何学习一个 Pareto 策略集，而不是单一策略。

也就是，你研究的不是“一个最优蓝方策略”，而是：

一组覆盖不同安全-业务-成本权衡的非支配防御策略。

子问题 3：怎样高效扩展 Pareto front

原始两阶段 C-MORL 已经提出了一个基本思路，但你现在如果以 AdaCS-DCS 为主，而且在 CybORG++ 上也有效，那么第三层科学问题就是：

如何在复杂 cyber defense 环境中，更有效地选择待扩展父策略，并动态调节约束强度，从而更高效地扩展 Pareto front。

\subsection{贡献解析}

贡献 1：问题建模贡献

我们将自主网络防御系统化建模为一个部分可观测、序贯对抗、多目标决策问题。

具体体现在：

把蓝方防御目标从单一安全收益扩展为 security / business / cost
强调偏好变化下的策略选择需求
将防御问题从“单策略最优”转化为“Pareto 策略集发现与分配”

这条贡献回答的是：
你到底在研究什么问题。

贡献 2：方法贡献

我们提出了一个面向网络防御的 AdaCS-DCS 增强型两阶段 Pareto 策略发现方法。

这条可以写成你的核心算法贡献。

如果你实验真的已经证明它比原始 Stage-2 更稳定、更强，那么可以写：

AdaCS：自适应地选择更值得扩展的父策略
DCS：针对不同策略、目标方向与扩展阶段动态调节约束强度
二者结合后，提升了 Pareto front 的扩展效率与质量

这里可以把名字写成：

AdaCS-DCS-CMORL
或者
an adaptive constrained Pareto extension framework for cyber defense

这条贡献回答的是：
你提出了什么新方法。

贡献 3：系统迁移与环境贡献

我们将两阶段多目标 Pareto 学习框架从简化环境推进到更复杂的 CybORG++ 网络防御环境，并构建了适用于该环境的训练、评估与策略分配流程。

这条很重要，因为它会把你的工作从“toy setting 成功”抬到“更真实 cyber range 上成功”。

你可以强调：

从 MiniCAGE 到 CybORG++ 的迁移
在更复杂、更接近真实网络防御的环境中验证方法
建立统一的多目标训练与评估协议

这条贡献回答的是：
你的方法是否真的能用于更强场景。

贡献 4：实验与机制贡献

我们通过系统实验表明，所提出的方法在 Pareto front 质量、偏好适应能力以及关键网络安全语义指标上优于现有基线，并揭示了其有效性的机制。

这条里可以放：

比 Stage-1 更好
比原始 Stage-2 更好
比 weighted-sum / single-objective / heuristic baselines 更好
在 HV、EU、SP、Pareto Count 上更优
在关键资产保护、业务扰动、防御成本等语义指标上更合理

如果你还有消融，就更强：

AdaCS 的贡献
DCS 的贡献
二者组合的贡献


\subsection{对应的现实场景和问题}
现实场景：
本文关注的现实网络安全场景是企业或关键基础设施网络中的持续入侵处置问题。在这类场景中，攻击者在获得初始立足点后，往往会继续进行内部侦察、横向移动与关键资产逼近；与此同时，蓝队或安全运营中心并不能在完全观测的条件下直接掌握全局状态，而必须依据有限告警、主机状态和网络迹象，持续决定是否执行监测分析、主机隔离、恶意清除、系统恢复、诱骗部署等防御动作。现实中的处置目标也并非单一的“最大化拦截攻击”，而是要同时兼顾安全遏制效果、业务连续性与响应处置成本，因此该类任务天然具有对抗性、序贯性与多目标权衡特征。

问题定义：
基于上述现实需求，本文将自主网络防御建模为一个部分可观测、对抗性的多目标序贯决策问题。具体而言，蓝方智能体需要在攻击持续演化的过程中，根据局部观测选择防御动作，以同时优化三个相互冲突但又必须共同考虑的目标：其一是安全目标，即尽可能降低攻击扩散、主机失陷和关键资产受损风险；其二是业务目标，即尽可能减少隔离、恢复、误处置等操作对业务可用性和服务连续性的影响；其三是成本目标，即尽可能控制高强度响应、人工介入和运维执行所带来的资源消耗。由此，本文所研究的核心问题可以表述为：在部分可观测的红蓝对抗环境中，如何学习一组能够覆盖不同安全—业务—成本偏好的高质量防御策略，并在运行时根据实际偏好自适应地选择合适策略，以支持现实网络中的入侵遏制与业务连续性平衡决策


\section{博弈建模}

\subsection{将红蓝攻防建模为部分可观测随机博弈}

现实网络安全中的自主防御并不是孤立的单智能体控制问题，而是攻击者与防御者围绕网络资产、系统权限与业务可用性展开的持续对抗过程。攻击者试图通过侦察、利用、横向移动与影响执行等行为扩大失陷范围并逼近关键资产；防御者则需要通过监测分析、主机隔离、恶意清除、系统恢复与诱骗部署等手段抑制攻击进展。在这一过程中，双方行为共同作用于网络状态的演化，因此本文首先将完整攻防过程建模为一个\textbf{部分可观测随机博弈}（partially observable stochastic game, POSG）。

形式上，完整红蓝交互可表示为
\begin{equation}
\mathcal{G}
=
\left\langle
\mathcal{I},
\mathcal{S},
\mathcal{A}^{B},
\mathcal{A}^{R},
\mathcal{O}^{B},
\mathcal{O}^{R},
\mathcal{T},
\mathcal{Z},
\mathbf{r}^{B},
r^{R},
\gamma
\right\rangle,
\end{equation}
其中，$\mathcal{I}=\{B,R\}$ 表示蓝方（defender）与红方（attacker）；$\mathcal{S}$ 表示网络系统的真实安全状态空间；$\mathcal{A}^{B}$ 与 $\mathcal{A}^{R}$ 分别表示蓝方和红方动作空间；$\mathcal{O}^{B}$ 与 $\mathcal{O}^{R}$ 分别表示双方的观测空间；$\mathcal{T}(s_{t+1}\mid s_t,a_t^B,a_t^R)$ 表示联合动作驱动的状态转移概率；$\mathcal{Z}(o_{t+1}^B,o_{t+1}^R\mid s_{t+1},a_t^B,a_t^R)$ 表示观测生成过程；$\mathbf{r}^{B}$ 与 $r^{R}$ 分别表示蓝方与红方收益函数；$\gamma\in(0,1)$ 为折扣因子。

上述建模有两个核心特点。其一，网络安全攻防具有明显的\textbf{时序性}：当前动作不仅影响当下收益，还会改变未来攻击路径、防御可用性与关键资产暴露程度。其二，该问题具有显著的\textbf{部分可观测性}：蓝方通常只能通过告警、日志、主机状态摘要与流量线索间接推断攻击态势，而无法直接获得全局真实状态。因此，与将网络防御建模为静态优化或完全可观测控制不同，POSG 能更准确地描述现实中的红蓝交互本质。

\subsection{网络防御中的状态、动作与观测设计}

在上述 POSG 框架下，状态 $s_t\in\mathcal{S}$ 用于刻画时刻 $t$ 的真实网络安全态势。具体而言，$s_t$ 不仅包含各主机的安全状态，还包含与防御决策直接相关的业务与运维信息。本文关注的状态因素主要包括：主机是否失陷、权限等级是否提升、关键资产是否暴露或遭受影响、服务是否处于可用状态，以及当前防御资源或系统恢复状态等。由于网络中的攻击传播与防御处置通常具有路径依赖性，这些状态量共同决定了后续攻防行为的效果。

红蓝双方的动作空间体现了攻防能力的差异性。对红方而言，动作 $a_t^R\in\mathcal{A}^R$ 可抽象为侦察扫描、漏洞利用、横向移动、权限提升、持久化维持与影响执行等典型攻击行为；对蓝方而言，动作 $a_t^B\in\mathcal{A}^B$ 则对应于监测分析、主机隔离、账户或恶意进程清除、系统恢复、诱骗部署及其他防御操作。红蓝双方动作在每个时刻共同作用于系统演化，并通过状态转移函数 $\mathcal{T}$ 影响未来网络态势。

与经典完全可观测控制问题不同，蓝方并不能直接访问 $s_t$，而只能获得局部观测 $o_t^B\in\mathcal{O}^B$。该观测通常由告警信息、主机状态摘要、服务异常迹象以及系统记录等组成，本质上是不完整且带噪的。因此，蓝方策略一般应依赖于历史交互而非单步观测，即
\begin{equation}
\pi^B(a_t^B \mid h_t^B),
\qquad
h_t^B=(o_0^B,a_0^B,\dots,o_t^B),
\end{equation}
其中 $h_t^B$ 表示蓝方截至时刻 $t$ 的动作--观测历史。该定义反映出：在现实网络防御中，决策并非基于“已知全局状态”的静态选择，而是基于有限信息不断更新态势判断后的序贯响应。

\subsection{防御者的向量化效用建模}

与传统网络防御强化学习只优化单一安全收益不同，现实蓝队处置往往需要在多个相互冲突的目标之间做出权衡。若一味追求最强隔离和恢复，虽然可能提升安全效果，却会显著放大业务扰动和运维成本；反之，若过度追求低成本与低扰动，则可能放任攻击扩散并增加关键资产受损风险。基于这一现实特征，本文将蓝方效用建模为一个\textbf{向量值回报}，而不是单一标量奖励。

具体地，蓝方在时刻 $t$ 的即时回报定义为
\begin{equation}
\mathbf{r}_t^B
=
\left(
r_t^{\mathrm{sec}},
r_t^{\mathrm{biz}},
r_t^{\mathrm{cost}}
\right),
\end{equation}
其中 $r_t^{\mathrm{sec}}$ 表示安全目标收益，用于刻画失陷主机减少、关键资产保护、攻击扩散抑制等安全效果；$r_t^{\mathrm{biz}}$ 表示业务目标收益，用于刻画服务可用性保持、业务中断最小化以及低扰动处置；$r_t^{\mathrm{cost}}$ 表示成本目标收益，用于刻画处置动作、恢复开销和持续运维负担。本文统一采用“目标值越大越优”的约定，因此即使某些项在具体实现中表现为负代价，其优化方向仍然保持一致。

在此基础上，蓝方策略 $\pi^B$ 的长期效用可写为折扣向量回报
\begin{equation}
\mathbf{J}(\pi^B)
=
\mathbb{E}\left[
\sum_{t=0}^{\infty}\gamma^t \mathbf{r}_t^B
\right].
\end{equation}
由于三个目标通常彼此冲突，本文不假设存在一个天然且固定的标量化权重可以一劳永逸地表示所有运维偏好。相反，我们认为更合理的目标应当是学习一组覆盖不同权衡关系的非支配防御策略，为后续按偏好自适应分配提供基础。由此，网络防御问题在本文中被刻画为一个\textbf{多目标而非单目标}的对抗决策问题。

\subsection{从完整攻防博弈到固定攻击者下的防御者多目标决策问题}

尽管完整红蓝交互更准确地对应于 POSG，但直接在双边部分可观测博弈上进行共同学习会面临显著挑战。一方面，攻击者与防御者策略同时更新会引入强烈的非平稳性；另一方面，在网络安全环境中，局部观测、长时序依赖和高维状态空间会进一步放大训练难度。考虑到本文的核心目标是研究\textbf{蓝方在安全、业务与成本之间的多目标权衡机制}，而不是首先解决双边学习的稳定性问题，因此我们对完整博弈进行了可训练化简。

具体而言，本文采用\textbf{固定红方策略}的设定，将攻击者视为环境中的对抗性动态来源。在给定红方策略 $\pi^R$ 后，原始双边博弈可约化为蓝方单边优化问题：
\begin{equation}
\mathcal{M}_{\pi^R}
=
\left\langle
\mathcal{S},
\mathcal{A}^B,
\mathcal{O}^B,
\mathcal{T}_{\pi^R},
\mathcal{Z}_{\pi^R},
\mathbf{r}_{\pi^R}^{B},
\gamma
\right\rangle,
\end{equation}
其中，$\mathcal{T}_{\pi^R}$ 和 $\mathbf{r}_{\pi^R}^{B}$ 分别表示在红方固定策略诱导下的状态转移与蓝方向量奖励。形式上，蓝方所面对的转移动力学可写为
\begin{equation}
\mathcal{T}_{\pi^R}(s_{t+1}\mid s_t,a_t^B)
=
\sum_{a_t^R\in\mathcal{A}^R}
\pi^R(a_t^R\mid h_t^R)\,
\mathcal{T}(s_{t+1}\mid s_t,a_t^B,a_t^R).
\end{equation}

由于蓝方仍然只能访问局部观测 $o_t^B$，上述化简后的问题在严格意义上应被视为一个\textbf{多目标部分可观测马尔可夫决策过程}。若在具体实现中通过环境包装器为蓝方提供了满足马尔可夫性的状态表示，则该问题可进一步近似为多目标马尔可夫决策过程；但从问题本质出发，本文更强调其“固定攻击者条件下的防御者侧多目标部分可观测决策问题”属性。

这一化简并不否认完整攻防博弈的存在，而是将本文的研究重点明确限定为：\textbf{在固定攻击行为模式下，如何为防御者学习一组高质量 Pareto 防御策略，以适应不同的安全--业务--成本偏好。}这一设定既保留了网络安全环境中的对抗性与不确定性，又使后续两阶段 Pareto 策略发现与自适应策略分配成为可能。


\subsection{建模选择讨论：POSG、Stackelberg 与 MDP 化简的关系}

需要说明的是，本文并不将完整红蓝交互直接表述为经典的 Stackelberg 博弈。其原因在于，Stackelberg 博弈通常强调领导者与跟随者之间的一次性承诺与响应关系，而本文关注的网络攻防过程是一个随时间演化的、带有状态转移和局部观测的不确定动态系统。攻击者与防御者的决策不是单轮静态选择，而是在长期交互中不断影响后续状态、观测与收益。因此，从问题本质出发，将其建模为部分可观测随机博弈更为准确。

尽管如此，在直觉层面，本文所采用的固定攻击者设定可以被理解为一种带有 ``leader--follower'' 色彩的近似化简：攻击者策略先验给定，防御者在该攻击行为模式下学习最优响应。然而，这种解释仅用于说明防御者侧学习问题的策略意义，而不是本文的正式数学建模。就严格表述而言，完整问题首先是一个向量值的部分可观测随机博弈；在固定攻击者策略后，才进一步化简为防御者侧的多目标决策问题。

进一步地，本文也不直接将完整问题写成标准 MDP。只有在攻击者策略固定、且防御者能够获得满足马尔可夫性的状态表示时，化简后的单边问题才可写成多目标马尔可夫决策过程。若防御者仅依赖局部告警和历史观测进行决策，则更准确的表述应是多目标部分可观测马尔可夫决策过程。基于这一点，本文采用如下分层表述：完整红蓝系统建模为部分可观测随机博弈，而本文实际求解的是固定攻击者条件下的防御者侧多目标决策问题。这样的建模选择既保留了网络安全场景中的对抗本质，又使后续的策略学习与算法设计具有可训练性和可实现性。

\subsection{Pareto 最优防御策略集}

在固定攻击者策略 $\pi^R$ 的条件下，本文的目标不是学习一个唯一的标量最优防御策略，而是学习一组能够覆盖不同安全--业务--成本权衡关系的高质量非支配防御策略。为此，设蓝方策略为 $\pi^B$，其对应的折扣向量回报定义为
\begin{equation}
\mathbf{J}(\pi^B)
=
\mathbb{E}_{\pi^B,\pi^R}
\left[
\sum_{t=0}^{\infty}
\gamma^t \mathbf{r}_t^B
\right]
\in \mathbb{R}^m,
\end{equation}
其中 $m$ 表示目标维度；在本文中，$m=3$，分别对应安全目标、业务目标和成本目标。

基于向量回报，可以定义策略之间的 Pareto 支配关系。对于任意两个蓝方策略 $\pi_1^B$ 和 $\pi_2^B$，若满足
\begin{equation}
\mathbf{J}(\pi_1^B) \succeq \mathbf{J}(\pi_2^B),
\end{equation}
即对任意目标维度 $k$ 均有 $J_k(\pi_1^B) \geq J_k(\pi_2^B)$，并且至少存在一个维度 $k'$ 使得
\begin{equation}
J_{k'}(\pi_1^B) > J_{k'}(\pi_2^B),
\end{equation}
则称策略 $\pi_1^B$ Pareto 支配策略 $\pi_2^B$。若一个策略不被任何其他策略支配，则称其为非支配策略。所有非支配策略构成的集合记为 Pareto 最优策略集，记作
\begin{equation}
\Pi^\star = \{\pi_1^B,\pi_2^B,\dots,\pi_K^B\}.
\end{equation}

本文采用 Pareto 策略集而非单一策略作为求解目标，主要基于如下考虑。首先，网络防御中的运维偏好并非固定不变：在高风险时段，决策者可能更强调安全收益；在业务高峰期，决策者可能更重视低扰动和服务连续性；在资源受限情形下，响应成本又可能成为更重要的约束。其次，不同偏好下的最优防御动作往往并不一致，单一策略难以同时在所有偏好上保持最优。因而，与其学习一个依赖固定标量化权重的单策略解，不如学习一个覆盖不同权衡关系的 Pareto 防御策略集，再根据实际偏好进行自适应调用。本文后续的两阶段学习框架正是围绕这一目标展开：先构建初始 Pareto 候选策略集，再通过定向扩展不断提升其覆盖范围与前沿质量。

\subsection{偏好驱动的策略分配}

在获得 Pareto 最优防御策略集 $\Pi^\star$ 之后，仍需回答一个关键问题：在具体部署时，系统应当如何根据当前运维需求选择合适的防御策略。为此，本文引入偏好向量
\begin{equation}
\mathbf{w} = (w_1,\dots,w_m)\in \Delta^{m-1},
\end{equation}
其中 $\Delta^{m-1}$ 表示 $(m-1)$ 维概率单纯形，即满足 $w_k \ge 0$ 且 $\sum_{k=1}^{m} w_k = 1$。在本文中，$\mathbf{w}$ 用于表示当前决策者对安全、业务与成本三个目标的重要性分配。

给定偏好向量 $\mathbf{w}$ 后，本文通过效用函数在 Pareto 策略集中进行策略分配。最基本的形式为线性效用
\begin{equation}
U(\mathbf{w}, \mathbf{J}(\pi^B))
=
\mathbf{w}^{\top}\mathbf{J}(\pi^B).
\end{equation}
据此，可定义给定偏好下的策略选择规则为
\begin{equation}
\pi^B(\mathbf{w})
=
\arg\max_{\pi^B \in \Pi^\star}
U(\mathbf{w}, \mathbf{J}(\pi^B)).
\end{equation}
该规则的含义是：在离线阶段，系统先学习一组互不支配的候选防御策略；在在线部署阶段，不再重新训练，而是根据当前偏好直接从策略集中选择效用最大的策略执行。

这种 ``先发现策略集、后按偏好分配'' 的机制比单一固定策略更符合现实网络防御需求。一方面，不同组织、不同网络业务阶段对安全性、业务连续性与运维成本的容忍度往往不同；另一方面，运维偏好也可能随着威胁等级、关键业务时段与资源状态发生变化。因而，本文不将偏好仅视为训练时的辅助参数，而是将其作为部署时决定策略调用的重要输入。由此，Pareto 策略集与偏好驱动分配共同构成了一个能够支持动态防御决策的统一框架。

\subsection{与 AdaCS-DCS 的联系}

上述建模为本文后续方法设计提供了直接依据。在固定攻击者条件下，蓝方所面对的是一个多目标决策问题，其核心挑战不再是寻找单一最优策略，而是高效发现并扩展一个高质量的 Pareto 防御策略集。在这一背景下，本文的方法由两个连续阶段组成：第一阶段用于生成初始 Pareto 候选策略集，第二阶段用于在现有非支配策略基础上沿不同目标方向进行进一步扩展。

具体而言，第一阶段的目标是通过多偏好训练获得覆盖不同权衡关系的初始策略池，并将其中的非支配策略写入解集缓冲区。第二阶段则从当前 Pareto 前沿中选择若干父策略，在保持其他目标不过度退化的前提下，沿特定目标方向执行约束扩展，从而逐步提升前沿质量。本文提出的 AdaCS-DCS 正是在这一阶段发挥核心作用。

其中，AdaCS（Adaptive Candidate Selection）的作用在于回答 ``扩展谁'' 的问题。与仅基于稀疏度或 crowding distance 选择父策略不同，AdaCS 希望综合考虑候选策略的前沿覆盖价值、扩展潜力与约束风险，以更有针对性地挑选值得继续扩展的父策略。DCS（Dynamic Constraint Scheduling）的作用在于回答 ``如何扩展'' 的问题。由于不同父策略、不同目标方向以及不同扩展轮次对应的可行性边界并不一致，采用统一固定的约束强度往往难以兼顾探索效率与约束满足性。因此，DCS 通过对约束强度进行动态调节，使扩展过程能够根据当前路径特征自适应地控制推进幅度。

从这个意义上说，AdaCS-DCS 并不是脱离前述建模而独立存在的工程技巧，而是对 ``固定攻击者条件下如何更高效地发现 Pareto 最优防御策略集'' 这一核心问题的直接算法回答。前述博弈建模说明了为何本文应以防御者侧的多目标策略集学习为目标；而 AdaCS-DCS 则进一步回答了，在这样的目标之下，如何更有效地完成 Pareto 前沿的构建、扩展与部署。


\section{理论说明}
\subsection{网络防御问题的多目标本质}

与传统单目标控制任务不同，现实网络防御中的决策目标通常并不唯一。对于蓝方而言，防御动作的价值不仅体现在是否能够抑制攻击扩散、保护关键资产和降低失陷范围，还体现在其对业务连续性与运维资源消耗的影响上。换言之，同一防御动作往往同时影响安全收益、业务扰动和处置成本三个维度，而这三个维度之间通常存在天然冲突关系。

具体而言，采取更激进的隔离、恢复或清除措施，往往有助于提升安全效果，但也可能导致服务中断、业务降级或额外的人工处置开销；相反，若蓝方过度追求低成本或低扰动，则可能放松对攻击扩散和关键资产风险的控制。因而，网络防御中的优良策略并不对应于某个单一维度上的极致最优，而对应于多个目标之间的合理权衡。本文将这种权衡形式化为向量值回报
\begin{equation}
\mathbf{r}_t = \left(r_t^{sec},\, r_t^{biz},\, r_t^{cost}\right),
\end{equation}
并进一步考虑其长期折扣回报
\begin{equation}
\mathbf{J}(\pi)
=
\mathbb{E}_{\pi}
\left[
\sum_{t=0}^{\infty} \gamma^t \mathbf{r}_t
\right].
\end{equation}

上述形式化表明，本文所研究的问题不是 ``在一个给定目标上寻找最优防御策略''，而是 ``在多个相互冲突的目标之间寻找高质量权衡策略''。因此，从问题本质上看，自主网络防御更适合被建模为一个多目标决策问题，而不是单目标优化问题。只有在这一前提下，后续关于 Pareto 前沿、策略分配以及偏好自适应调用的讨论才具有理论上的一致性。

\subsection{单一标量化策略的局限性}

在多目标强化学习中，一种常见处理方式是将向量值回报通过固定权重进行标量化，从而将问题转化为单目标优化。例如，给定偏好向量
\begin{equation}
\mathbf{w} \in \Delta^{m-1},
\end{equation}
可定义标量化效用为
\begin{equation}
U(\mathbf{w}, \mathbf{J}(\pi))
=
\mathbf{w}^{\top}\mathbf{J}(\pi).
\end{equation}
在该设定下，训练过程本质上是在求解
\begin{equation}
\pi_{\mathbf{w}}^{\star}
=
\arg\max_{\pi} U(\mathbf{w}, \mathbf{J}(\pi)).
\end{equation}
这意味着，固定标量化策略只对应于某一个特定偏好向量 $\mathbf{w}$ 下的最优解。

然而，现实网络防御中的偏好通常不是静态不变的。不同组织对业务中断的容忍度不同，不同业务阶段对服务可用性的要求不同，不同威胁等级下对安全风险的容忍度也不同。因而，一个在偏好 $\mathbf{w}_1$ 下训练得到的最优策略，并不能保证在另一组偏好 $\mathbf{w}_2$ 下仍然最优。特别地，当两个偏好向量强调的目标维度不同，所诱导出的最优防御策略也可能完全不同。由此可见，单一固定权重下的标量化训练只能覆盖多目标问题中的局部决策需求，而无法普遍应对偏好变化场景。

从决策角度看，单一策略的局限性在于它将多目标问题提前压缩为某一种固定价值取向，从而牺牲了运行时根据具体任务需求调整策略的能力。对于网络防御而言，这种局限尤其明显：当系统处于高风险时段时，决策者可能更愿意牺牲业务和成本以保障安全；而在业务高峰或资源受限阶段，决策者又可能更强调低扰动和低开销。因此，仅依赖一个固定标量化策略，不足以支撑现实网络中的动态防御需求。

\subsection{Pareto 防御策略集的必要性}

既然单一固定偏好策略无法充分覆盖现实中的多样化防御需求，更自然的思路便是学习一个能够覆盖不同权衡关系的策略集合。本文采用的核心思想是，不直接寻找唯一最优策略，而是寻找一组在向量回报空间中互不支配的防御策略，即 Pareto 防御策略集。

设蓝方策略 $\pi_1$ 和 $\pi_2$ 的长期向量回报分别为 $\mathbf{J}(\pi_1)$ 和 $\mathbf{J}(\pi_2)$。若满足
\begin{equation}
\forall k,\; J_k(\pi_1) \ge J_k(\pi_2),
\end{equation}
并且至少存在一个维度 $k'$ 使得
\begin{equation}
J_{k'}(\pi_1) > J_{k'}(\pi_2),
\end{equation}
则称 $\pi_1$ 支配 $\pi_2$。所有不被其他策略支配的策略构成 Pareto 策略集，记作
\begin{equation}
\Pi^{\star} = \{\pi_1,\pi_2,\dots,\pi_K\}.
\end{equation}

Pareto 策略集的价值在于，它保留了不同目标权衡下的候选解结构，使系统能够在部署阶段根据当前偏好进行灵活选择，而不必在训练阶段预先锁定唯一权重。从理论上看，若偏好会变化，则不存在一个单一策略能够在所有偏好下同时最优；而 Pareto 策略集恰恰提供了一组候选非支配解，使得不同偏好都能在其中找到对应的高质量策略。对于网络防御这一类目标冲突显著、偏好动态变化的问题而言，这种 ``先学习策略集，再按偏好调用'' 的范式比 ``先固定偏好，再学习单策略'' 更具一般性。

从现实意义上看，Pareto 防御策略集也更符合安全运营的实际模式。组织决策者通常并不希望系统只输出一种刚性的统一防御方案，而更希望系统能够根据威胁等级、业务重要性、服务窗口和资源状态，选择更偏安全、更偏业务连续性或更偏低成本的处置策略。因此，学习 Pareto 防御策略集并在运行时执行偏好驱动分配，不仅具有数学上的合理性，也更符合真实网络防御的部署需求。

\subsection{两阶段 Pareto 策略发现框架的合理性}

在明确目标是学习 Pareto 防御策略集之后，进一步需要回答的问题是：如何在复杂的网络防御环境中高效地构造该策略集。直接从零开始同时逼近整个 Pareto 前沿通常较为困难，其原因在于：一方面，多目标策略空间本身具有较高复杂度；另一方面，网络防御环境中的状态转移、对抗行为和部分可观测性会进一步增加搜索难度。基于此，本文采用两阶段策略发现框架，将 Pareto 前沿构建问题分解为 ``先覆盖、后扩展'' 两个相对可控的子问题。

第一阶段的目标是获得一个能够覆盖不同偏好方向的初始策略集。具体而言，通过对多个偏好向量分别训练防御策略，可以在目标空间中快速得到若干具有代表性的候选解。这一阶段不要求一次性逼近完整 Pareto 前沿，而是强调形成一个足够好的初始非支配解集，为后续扩展提供起点。换言之，第一阶段承担的是 ``建立前沿骨架'' 的作用。

第二阶段的目标是在已有候选解基础上进一步提升 Pareto 前沿质量。与第一阶段的全局性初始探索不同，第二阶段从当前非支配策略出发，沿特定目标方向执行受约束扩展，从而逐步发现新的高质量非支配点。这样的设计具有明显的理论直觉：当初始策略集已经提供了若干可行权衡点后，更高效的做法不是重新在整个策略空间中盲目搜索，而是围绕现有前沿执行定向改进。通过这一过程，策略集可以在保持已有解结构的同时，逐步提升覆盖范围、前沿密度和整体质量。

因此，两阶段框架的合理性体现在两个方面。其一，它将复杂的 Pareto 前沿构建任务分解为初始覆盖与定向扩展两个子过程，降低了直接求解的难度；其二，它与网络防御问题的实际需求相契合，即先快速形成一组可用防御策略，再围绕这些策略持续优化其安全--业务--成本权衡表现。基于这一思路，本文后续的方法设计不再将策略学习视为单点最优搜索，而将其视为面向 Pareto 防御策略集的分阶段发现过程。


\subsection{受约束扩展机制的理论解释}

在多目标防御策略发现中，第二阶段的核心任务不是无约束地优化某一个单独目标，而是在推动某一目标进一步改善的同时，显式控制其他目标不要发生不可接受的退化。设当前父策略为 $\pi$，其向量回报为 $\mathbf{J}(\pi)$，若目标扩展方向为第 $j$ 个目标，则第二阶段的理想优化问题可写为
\begin{equation}
\max_{\pi'} \; J_j(\pi')
\quad
\text{s.t.}
\quad
J_k(\pi') \ge \beta J_k(\pi) - \varepsilon_k,\;\; \forall k \neq j,
\end{equation}
其中 $\beta$ 为约束强度参数，$\varepsilon_k$ 为允许的容差项。上述形式表明，阶段二本质上是在局部前沿上执行一种 ``受控改进''：系统并不允许为了提升某一个目标而无节制地牺牲其他目标，而是要求候选子策略在非目标维度上仍然维持一定质量。

这种约束扩展机制与现实网络防御中的决策逻辑高度一致。对于蓝方而言，提升安全效果固然重要，但若这种提升伴随着过高的业务中断或处置代价，则该改进在实际场景中往往不可接受；反之，若一味压低成本而显著削弱安全防御能力，该策略同样缺乏部署价值。因此，阶段二并不是简单地沿某个目标方向做单目标爬升，而是在多个目标之间执行局部可行的折中推进。换言之，其理论意义不在于追求单个目标的极端最优，而在于推动策略集向更优的 Pareto 区域演化。

从 Pareto 几何角度看，阶段一得到的初始非支配策略集提供了若干离散的前沿支撑点，而阶段二的受约束扩展则试图在这些已有支撑点附近进一步向外推进。若某个候选子策略在目标维度 $j$ 上取得提升，同时在其他维度上的退化受到约束控制，则该子策略至少有可能成为新的非支配候选点，从而使前沿向外扩张。由此，阶段二的受约束优化可以被理解为一种面向 Pareto 前沿 ``局部外推'' 的机制：它不保证每次扩展都得到新的非支配点，但它保证了扩展方向始终围绕 ``改进一个目标且不显著破坏其余目标'' 这一原则展开。

进一步地，这一机制也解释了为何本文将阶段二设计为沿不同目标方向分别执行扩展。由于不同目标对应不同的实际防御偏好，若仅围绕单一方向进行改进，所得策略集很容易偏向某一类价值取向；而通过对每个目标方向分别执行受约束扩展，可以系统地提升整个策略集对不同偏好区域的覆盖能力。因此，受约束扩展并非阶段一训练的简单重复，而是在已有 Pareto 结构基础上执行方向化、局部化和可控化的前沿推进过程。

\subsection{AdaCS 的理论动机}

在两阶段 Pareto 策略发现框架中，阶段二的效果不仅取决于扩展规则本身，还取决于 ``从哪些父策略出发进行扩展''。若父策略选择不当，即使后续扩展机制合理，也可能因为起点过于拥挤、扩展空间有限或约束风险过高，而无法产生高质量的新非支配点。基于这一点，本文认为父策略选择本身应被视为一个独立且关键的问题，而不应被简化为仅依据几何稀疏度进行机械筛选。

传统做法常采用 crowding distance 等几何指标来选择扩展父策略，其直觉在于：位于稀疏区域的策略更有可能补齐前沿空白。然而，仅依赖几何稀疏性是不充分的。首先，一个几何上稀疏的点并不必然具有较大的扩展潜力；若该点已经接近局部可行边界，则继续扩展可能很快失败。其次，某些策略虽然位于相对拥挤区域，但若其在特定目标方向上仍有较大可提升空间，且约束满足风险较低，则它依然可能产生更优的新候选解。最后，从部署角度看，真正有价值的新策略不应只改善局部几何分布，还应扩大对不同偏好区域的覆盖能力。

基于上述考虑，AdaCS 的核心思想是将 ``值得扩展'' 从单一几何概念提升为综合的扩展价值评估。设父策略为 $\pi_i$，其选择得分可抽象表示为
\begin{equation}
S_{\mathrm{AdaCS}}(\pi_i)
=
\lambda_1 C_i
+
\lambda_2 P_i
-
\lambda_3 R_i
+
\lambda_4 G_i,
\end{equation}
其中，$C_i$ 表示几何稀疏性或 crowding 信息，$P_i$ 表示扩展潜力，$R_i$ 表示约束失败风险，$G_i$ 表示对偏好覆盖的潜在增益，$\lambda_1,\lambda_2,\lambda_3,\lambda_4$ 为相应权重。该表达式并不依赖某一种特定实现，而是强调：父策略选择应同时考虑其在目标空间中的位置、进一步推进的可能性以及生成新非支配点的价值。

从理论解释上看，AdaCS 的必要性在于它使阶段二的搜索从 ``被动补点'' 转向 ``主动发现有价值的扩展路径''。如果某一父策略同时具备较大的扩展潜力、较低的约束风险以及更高的偏好覆盖价值，那么优先从该策略出发扩展，更有可能提高单位计算预算下发现新非支配点的效率。换言之，AdaCS 的目标并不是单纯追求前沿的局部均匀性，而是通过对扩展价值进行显式建模，使有限的扩展次数能够更集中地投入到最有希望改善 Pareto 前沿质量的路径上。

因此，AdaCS 的理论动机可以概括为：在多目标防御场景中，父策略选择不应仅由当前几何结构决定，还应由策略的潜在改进空间、约束可行性以及部署层面的覆盖意义共同决定。只有在这一更丰富的选择准则下，阶段二才能真正发挥面向高质量 Pareto 前沿的定向扩展作用。

\subsection{DCS 的理论动机}

除父策略选择之外，阶段二的另一个关键问题是 ``以多严格的约束进行扩展''。若对所有父策略、所有目标方向和所有扩展轮次都采用同一固定约束强度，则隐含地假设：不同扩展路径具有相同的可行域形状、相同的风险水平以及相同的改进空间。然而，在复杂网络防御环境中，这一假设通常并不成立。不同父策略位于目标空间中的位置不同，所对应的安全--业务--成本平衡状态不同；不同扩展方向也意味着不同的改进难度与代价；而随着扩展轮次推进，局部可行域与前沿结构还会进一步变化。因此，单一固定约束强度往往难以同时兼顾可行性与推进幅度。

若约束过严，则候选子策略很难满足非目标维度的保持要求，导致扩展频繁失败，策略集难以向外推进；若约束过松，则虽然可行解更容易生成，但这些候选解可能通过过度牺牲其他目标来换取局部提升，从而削弱 Pareto 前沿的整体质量。由此可见，约束强度本质上控制着 ``保守性'' 与 ``探索性'' 之间的平衡：越保守，越容易保持现有目标平衡但难以扩展；越激进，越容易产生大幅推进但也更容易破坏非目标维度。

基于这一观察，DCS（Dynamic Constraint Scheduling）的核心思想是：约束强度不应是全局固定常数，而应根据父策略特征、扩展方向以及搜索进度进行动态调节。形式上，可将路径相关的约束强度写为
\begin{equation}
\beta_{i,j,r}
=
f\!\left(
\phi_i,\,
d_j,\,
q_r
\right),
\end{equation}
其中，$\phi_i$ 表示父策略 $\pi_i$ 的状态特征或前沿特征，$d_j$ 表示第 $j$ 个目标方向，$q_r$ 表示当前扩展轮次或进度信息。函数 $f(\cdot)$ 的具体实现可以有多种形式，但其理论含义是一致的：不同扩展路径应当拥有不同的约束调度策略。

从机制上看，DCS 的价值在于为不同路径提供差异化的推进节奏。对于约束边界较紧、失败风险较高的路径，较为宽松或更具缓冲的约束有助于维持可行性；对于扩展潜力较大、非目标维度更稳定的路径，则可以采用更积极的推进强度，以提高前沿外推幅度。随着搜索不断进行，前沿结构也会发生变化，因而约束调度还应具备阶段性：早期可以偏向更强探索以拓展覆盖范围，后期则可逐步强调更细致和更稳定的局部优化。

因此，DCS 的理论动机可以概括为：阶段二的约束设计不应被视为静态超参数选择问题，而应被视为与父策略、目标方向和扩展进程耦合的动态决策问题。动态约束调度的作用，不是简单替代固定参数，而是更系统地平衡 ``扩展成功率'' 与 ``前沿改进幅度'' 之间的内在张力，从而提升阶段二在复杂多目标防御环境中的适应性与效率。

\subsection{命题与讨论}

基于上述理论解释，本文进一步给出若干支撑性命题，用于说明 Pareto 策略集、受约束扩展以及 AdaCS-DCS 设计的必要性。这些命题不旨在构成完整的收敛理论，而是用于刻画本文方法设计背后的关键逻辑。

\paragraph{命题 1.}
若 Pareto 前沿中存在至少两个互不支配且回报向量不共线的策略 $\pi_1$ 与 $\pi_2$，则不存在单一策略 $\hat{\pi}$ 能够在所有偏好向量 $\mathbf{w}\in\Delta^{m-1}$ 下同时满足
\begin{equation}
U(\mathbf{w}, \mathbf{J}(\hat{\pi}))
=
\max_{\pi} U(\mathbf{w}, \mathbf{J}(\pi)).
\end{equation}

该命题表明，只要网络防御问题中确实存在多个互不支配的高质量权衡解，则单一固定策略不可能在所有偏好下同时保持最优。因而，学习 Pareto 防御策略集并在部署时进行偏好驱动分配，不是方法上的可选修饰，而是多目标问题本身所要求的自然结果。

\paragraph{命题 2.}
设父策略为 $\pi$，目标扩展方向为第 $j$ 个目标。若候选子策略 $\pi'$ 满足
\begin{equation}
J_j(\pi') > J_j(\pi),
\qquad
J_k(\pi') \ge \beta J_k(\pi) - \varepsilon_k,\;\; \forall k \neq j,
\end{equation}
则父策略 $\pi$ 不可能在所有目标维度上严格支配 $\pi'$。

该命题说明，受约束扩展所生成的可行子策略至少摆脱了 ``被父策略直接严格支配'' 的情形。虽然它仍可能被其他已有策略支配，但只要目标维度取得提升且非目标维度退化受到控制，该子策略就具备成为新 Pareto 候选点的基本可能性。这一性质为阶段二沿不同目标方向推进 Pareto 前沿提供了理论上的合理依据。

\paragraph{命题 3.}
若两个候选父策略 $\pi_a$ 和 $\pi_b$ 满足：$\pi_a$ 的预期扩展潜力不低于 $\pi_b$，约束失败风险不高于 $\pi_b$，且偏好覆盖增益不低于 $\pi_b$，则在相同扩展预算下，优先扩展 $\pi_a$ 比优先扩展 $\pi_b$ 更有可能生成高质量新候选解。

该命题并不声称存在统一闭式最优选择规则，而是强调：父策略选择本身具有显著的价值差异。也就是说，阶段二扩展效率不仅取决于 ``如何扩展''，还取决于 ``从哪里扩展''。这为 AdaCS 通过综合评估扩展价值来选择父策略提供了直接动机。

\paragraph{命题 4.}
若不同扩展路径的可行边界、目标推进空间与约束风险存在显著差异，则不存在单一固定约束强度可以在所有路径上同时达到最优的 ``可行性--推进幅度'' 平衡。

这一命题表明，固定约束强度在异质扩展路径下不可避免地会产生失配：对某些路径过严，对另一些路径过松。因而，约束调度应当依赖于路径上下文而非统一常数，这正是 DCS 设计的核心理论基础。

综上，上述命题共同支持了本文的整体方法逻辑：网络防御的多目标本质决定了需要 Pareto 策略集而非单一策略；Pareto 前沿的构建需要受约束扩展而非无约束单目标优化；而在受约束扩展阶段，父策略选择与约束强度调度又分别构成影响搜索效率与前沿质量的两个关键杠杆。AdaCS-DCS 正是在这一理论脉络下提出的，它并非对原有框架的简单工程调整，而是针对 ``如何更高效地发现高质量多目标防御策略集'' 这一核心问题给出的直接算法回答。



\section{文章升级方向-升级为攻击者不确定性下的鲁棒受限多目标强化学习}


将现有方法从单一脚本化攻击者条件下的蓝方多目标强化学习，升级为攻击者策略族不确定性下的鲁棒受限多目标强化学习（opponent-robust constrained multi-objective reinforcement learning, OR-CMORL），能够显著提升论文的创新性与研究完整性。但这种提升并不意味着研究会直接从“应用型迁移论文”跃升为“完整博弈均衡求解论文”，而是意味着本文的研究主线将从
\emph{cyber 场景中的 C-MORL 方法迁移与增强}
进一步提升为
\emph{攻击者不确定性下的 opponent-robust constrained MORL}。

原始 C-MORL 的公开主线是两阶段 Pareto front discovery，即首先初始化一组候选策略，再利用受约束扩展机制补全 Pareto front，其核心证据主要体现为 HV、EU、SP 等多目标评价指标上的实证提升\cite{CMORL_ICLR2025}。另一方面，CybORG 环境本身是为自主网络攻防研究设计的对抗式仿真平台\cite{CybORG2020}。因此，将现有方法从“单一 scripted red 下的 blue-only MORL”扩展为“攻击者策略族上的 robust C-MORL”，能够使论文在科学问题、方法定义以及理论分析层面都更加完整。

\subsection{对创新性的影响}

从创新性的角度看，这一修改将带来较为实质性的提升。在单一固定红方设定下，现有工作本质上更接近于以下三类贡献的组合：其一，将原方法迁移到新的 cyber 对抗场景；其二，对问题进行新的任务表述；其三，在工程实现上对原始方法进行增强。然而，当研究对象升级为攻击者策略族或多种红方策略下的鲁棒防御问题后，论文的创新结构将进一步扩展。

具体而言，第一，本文研究的理论对象将从普通多目标决策问题，提升为对手不确定性下的鲁棒多目标决策问题；第二，方法主张将从针对原始 C-MORL 的场景增强，进一步上升为 OR-CMORL，即面向攻击者不确定性的鲁棒受限多目标强化学习；第三，方法的部署意义也将从“对单一红方有效”提升为“对一类攻击行为均具有更强稳健性”。因此，论文的创新强度将不再局限于“方法迁移与实验增强”，而是进一步触及问题定义与方法机制本身的扩展。

当然，这种创新性的提升成立有一个前提，即训练目标、理论分析和实验验证三者必须保持一致。若仅仅是在评测阶段增加多个红方策略，而训练目标和方法机制仍然保持单红方设定，则创新性提升将较为有限。相反，只有当训练目标被明确改写为鲁棒形式，理论分析相应升级为鲁棒版本，并且实验能够证明方法在多红方条件下确实更稳定、更强时，这一修改才构成实质性的学术贡献。

\subsection{博弈建模与理论分析的变化}

这一修改会使本文的建模方式和理论分析都发生较大变化。原有工作通常采用两层结构来表述问题：首先，将红蓝攻防过程描述为一个完整的部分可观测随机博弈（POSG）；其次，在求解时将问题化简为固定红方条件下的 defender-side MO-POMDP 或 MO-MDP。此时，蓝方的任务本质上是针对一个固定攻击者策略求解多目标最优响应。

在引入鲁棒设定之后，更合适的表述方式应当扩展为三层结构。第一层仍然是完整红蓝对抗的向量值 POSG，用于刻画网络攻防的整体交互过程；第二层将研究对象从单一攻击者策略 $\nu$ 下的蓝方问题，升级为攻击者策略族
\[
\Pi_R=\{\nu_1,\nu_2,\dots,\nu_K\}
\]
下的蓝方鲁棒决策问题；第三层则是在训练阶段对这一鲁棒问题进行可计算化简，将蓝方在单一红方条件下的目标
\[
G_i(\pi_B,\nu)
\]
进一步推广为鲁棒目标，例如最坏情形定义
\[
J_i^{\mathrm{rob}}(\pi_B)=\min_{\nu\in\Pi_R} G_i(\pi_B,\nu),
\]
或者风险敏感形式
\[
J_i^{\mathrm{cvar}}(\pi_B)=\mathrm{CVaR}_{\alpha}\bigl(G_i(\pi_B,\nu)\bigr).
\]

在这样的建模下，本文不再是“针对某一个攻击脚本学习最优防御响应”，而是“针对一类攻击行为模式学习鲁棒 Pareto 最优防御策略集”。这种问题表述更加契合 CybORG 一类自主网络攻防环境的设计初衷，即蓝方在部署时需要面对多种潜在攻击策略，而非单一、固定且可预知的攻击脚本\cite{CybORG2020}。

与建模变化相对应，理论分析也需要整体升级为鲁棒版本。首先，原本讨论的普通向量回报函数 $\mathbf{J}(\pi)$，应当被鲁棒向量回报函数 $\mathbf{J}^{\mathrm{rob}}(\pi)$ 所替代。其次，原本基于单一环境回报定义的 Pareto front，需要被重新定义为基于鲁棒聚合回报的 robust Pareto front，即支配关系建立在最坏情形或风险敏感聚合后的性能向量之上。再次，原始 C-MORL 中 Stage-2 的受约束扩展机制，需要进一步改写为 \emph{robust constrained extension}，使其能够在攻击者策略族的不确定性下仍然有效地补全鲁棒 Pareto front\cite{CMORL_ICLR2025}。最后，现有方法中的 feasibility gate 也可以从经验性的工程筛选机制，进一步提升为具有统计含义的高概率可行性判别机制。例如，当无法对所有攻击者策略进行精确遍历评估时，可以基于有限红方集合的采样结果构造 lower-confidence bound，从而得到“以至少 $1-\delta$ 的概率满足真实鲁棒约束”的可行性保证。

因此，从理论视角看，这一修改的本质在于：本文的研究对象将从 fixed-opponent 条件下的普通 Pareto policy discovery，升级为 adversary-family uncertainty 条件下的 robust Pareto policy discovery。这一变化构成了论文最关键的学术提升。

\subsection{方法升级的可行性分析}

尽管上述修改幅度较大，但从实现角度看，这一路线仍然是可行的，而且明显比直接升级为 Nash 均衡求解或 Stackelberg 博弈求解更加现实。其原因在于，现有代码基础已经具备实现鲁棒版本所需的多个关键组件，包括三目标 cyber reward 设计、两阶段 C-MORL 主链路、IPO-style constrained extension，以及 feasibility gate 等机制。这意味着本文并不需要从零开始实现一个全新的鲁棒强化学习框架，而是在已有 C-MORL 迁移实现的基础上，将单红方目标定义替换为多红方聚合的鲁棒目标定义。

当然，这种升级仍然面临若干实际难点。第一，训练目标的改写并非简单替换即可完成。若直接采用
\[
\min_{k} G_i(\pi,\nu_k)
\]
作为训练信号，会带来不可导、方差增大以及训练不稳定等问题。因此，在训练阶段更现实的做法通常是采用 soft-min 近似，而在最终评估阶段再回到 hard min 或 CVaR 定义。第二，采样预算将显著增加。原本每个策略只需在单一红方下评估，现在则需要在 $K$ 个红方策略下进行评估，这将导致 Stage-1、Stage-2 以及最终 evaluation 的成本同步上升。因此，攻击者策略族的规模、每个红方的评估 episode 数量，以及 feasibility gate 的触发频率都需要谨慎控制。第三，理论与实验口径必须严格统一。若理论部分基于 worst-case 或 CVaR 定义鲁棒目标，而实验部分仅采用多红方平均回报进行汇报，则会导致整篇论文在问题定义和实验验证之间产生不一致。

基于上述考虑，一条更稳妥的实现路线是采用渐进式升级策略。首先，构造一个规模较小但具有代表性的红方策略族，例如包含 $3$ 至 $5$ 个具有不同攻击风格的代表性红方；其次，优先将评估阶段改写为鲁棒版本，以观察 single-red 训练策略在 multi-red worst-case 条件下的真实性能退化；然后，在此基础上将 feasibility gate 升级为 robust gate，使 Stage-2 在接受候选策略时考虑多红方条件下的鲁棒可行性；最后，再进一步将训练目标升级为 robust soft-min 版本。这一路线的优点在于可以先验证新问题设定是否成立，再逐步将方法机制提升到与理论主张一致的层次，而不必一次性重写全部训练框架。

\subsection{论文定位的变化}

如果上述升级能够顺利完成，那么本文最合适的论文定位将不再是“在 CybORG 场景中迁移和增强 C-MORL”，而应当被表述为：
\begin{quote}
Opponent-robust constrained multi-objective reinforcement learning for cyber defense under attacker uncertainty.
\end{quote}

这一定位强调，本文的核心贡献并不是求解完整的红蓝博弈均衡，而是在攻击者策略不确定性条件下，学习一组鲁棒的、多目标受约束防御策略。这样的定位既与现有代码基础和可实现方法保持一致，也能够在学术上体现出比单一红方设定更强的问题深度与方法价值。

\subsection{小结}

综上所述，将现有工作从单一 scripted red 条件下的 blue-only C-MORL，升级为攻击者策略族不确定性下的 opponent-robust constrained MORL，将显著提升论文的创新性、问题完整性与理论层次。其本质上的变化可以概括为：在建模上，从 fixed-opponent 的多目标最优响应问题，升级为 adversary-family uncertainty 下的鲁棒 Pareto policy discovery；在方法上，从普通两阶段 C-MORL 扩展为 robust C-MORL；在实验上，从对单一攻击者有效，提升为对多类攻击行为更稳定、更具泛化能力。与此同时，这一路线在现有代码和实验基础上具有较高可行性，并且明显比直接转向 Nash 或 Stackelberg 均衡求解更现实。因此，这是一条兼顾学术提升与工程可行性的合理升级路径。


\end{document}
